\documentclass[../mathNotesPreamble]{subfiles}

\providecommand{\relscalefact}{1.4}
\begin{document}
\relscale{\relscalefact}
  \section{10.2: Matrix Representations of Graphs}
  \begin{defn*}
    An $m\times n$ \textbf{matrix $A$ over a set $S$} is a rectangular array of elements of $S$ arranged into $m$ rows and $n$ columns:
      \[A=\begin{bmatrix}
        a_{11}& a_{12}& \dots& a_{1j}& \dots& a_{jn}\\
        a_{21}& a_{22}& \dots& a_{2j}& \dots& a_{2n}\\
        \vdots& \vdots& \ddots& \vdots& \ddots& \vdots\\
        a_{i1}& a_{i2}& \dots& a_{ij}& \dots& a_{in}\\
        \vdots& \vdots& \ddots& \vdots& \ddots& \vdots\\
        a_{m1}& a_{m2}& \dots& a_{mj}& \dots& a_{mn}\\
      \end{bmatrix}\]
  \end{defn*}
    \emph{Note}: Two matrices $A$ and $B$ are only equal if every component of $A$ and $B$ are equal:
      \[A=B \Leftrightarrow a_{ij}=b_{ij}\hspace*{8pt} \forall i=1,2,\dots,m, \textnormal{ and } \forall j=1,2,\dots,n.\]
  \vspace*{\stretch{1}}

  \begin{defn*}
    Let $G$ be a directed graph with ordered vertices $v_1, v_2,\dots, v_n$. The \textbf{adjacency matrix of $G$} is the $n\times n$ matrix $A=(a_{ij})$ over the set of nonnegative integers such that
      \[a_{ij}=\textnormal{ the number of arrows from } v_i \textnormal{ to } v_j,\ \forall i,j=1,2,\dots,n.\]
  \end{defn*}

  \noindent
  \begin{minipage}{0.475\linewidth}
    \begin{center}
      \begin{tikzpicture}[
        node distance=30mm,
        myNode/.style={circle, fill=black, inner sep=1.5pt},
        every loop/.style={looseness=30, ->},
        label distance=2pt
      ]
        \node[myNode, label=-90:$v_1$] (v1) at (0,0) {}
          edge [in=90,out=180,loop]
          node[above left] {$e_1$} (v1);
        \node[myNode, right=of v1, label=0:$v_2$]  (v2) {}
          edge [in=45,out=135,loop]
          node[above] {$e_3$} (v2);
        \node[myNode, below right=of v1, label=-180:$v_3$]  (v3) {};
%
        \draw[->] (v2) to [bend right=45] node[above left] {$e_4$} (v3);
        \draw[->] (v2) to [bend left=45] node[below right] {$e_5$} (v3);
        \draw[->] (v2) -- (v1) node[pos=0.5, above] {$e_2$};
        \draw[->] (v3) -- (v1) node[pos=0.5, below left] {$e_6$};
      \end{tikzpicture}
    \end{center}
  \end{minipage}\hspace*{\stretch{1}}%
  \begin{minipage}{0.5\linewidth}
    \begin{align*}
      A&=
        \begin{blockarray}{cccc}
          & v_1& v_2& v_3\\
          \begin{block}{c[ccc]}
            v_1& 1& 0& 0\\
            v_2& 1& 1& 2\\
            v_3& 1& 0& 0\\
          \end{block}
        \end{blockarray}
    \end{align*}
  \end{minipage}
  \pagebreak

  \begin{ex*}
    Let
      \[A=
      \begin{bmatrix}
        0& 1& 1& 0\\
        1& 1& 0& 2\\
        0& 0& 1& 1\\
        2& 1& 0& 0
      \end{bmatrix}
      \]
    Draw a directed graph that has $A$ as its adjacency matrix.
    \vspace*{\stretch{1}}
  \end{ex*}
  \begin{ex*}
    Find the adjacency matrix for the following graph:

    \begin{tikzpicture}[
      node distance=30mm,
      myNode/.style={circle, fill=black, inner sep=1.5pt},
      every loop/.style={looseness=40},
      label distance=2pt
    ]
      \node[myNode, label=180:$v_1$] (v1) at (0,0) {};
      \node[myNode, right=of v1, label=0:$v_2$]  (v2) {}
        edge [in=45,out=135,loop]
        node[above] {$e_3$} (v2);
      \node[myNode, below=of v2, label=0:$v_3$]  (v3) {};
      \node[myNode, below=of v1, label=-180:$v_4$]  (v4) {}
        edge [in=-45,out=-135,loop]
        node[below] {$e_6$} (v4);
%
      \draw (v1) -- (v4) node[pos=0.5, left] {$e_1$};
      \draw (v2) to [bend left=25] node[below right] {$e_4$} (v3);
      \draw (v2) to [bend right=25] node[below left] {$e_5$} (v3);
      \draw (v1) -- (v2) node[pos=0.5, above] {$e_2$};
      \draw (v2) -- (v4) node[pos=0.5, above left] {$e_7$};
    \end{tikzpicture}
  \end{ex*}
  \pagebreak

  \begin{ex*}
    Find the adjacency matrix of the following disconnected graph:

%    \noindent
%    \begin{minipage}{0.475\linewidth}
      \begin{center}
        \begin{tikzpicture}[
          node distance=25mm and 20mm,
          myNode/.style={circle, fill=black, inner sep=1.5pt},
          every loop/.style={looseness=30},
          label distance=2pt
        ]
          \node[myNode, label=-90:$v_1$] (v1) at (0,0) {}
            edge [in=90,out=180,loop]
            node[above left] {$e_1$} (v1);
          \node[myNode, right=of v1, label=0:$v_2$]  (v2) {};
          \node[myNode, below right=of v1, label=-180:$v_3$]  (v3) {};
          \node[myNode, right=of v2, label=0:$v_4$]  (v4) {};
          \node[myNode, right=of v3, label=0:$v_5$] (v5) {}
            edge [in=-45,out=-135,loop]
            node[below right] {$e_6$} (v5);
          \node[myNode, right=of v4, label=0:$v_6$]  (v6) {};
          \node[myNode, right=of v5, label=0:$v_7$]  (v7) {};
%
          \draw (v1) -- (v3) node[pos=0.5, below left] {$e_2$};
          \draw (v2) to [bend right=25] node[above left] {$e_3$} (v3);
          \draw (v2) to [bend left=25] node[below right] {$e_4$} (v3);
          \draw (v4) -- (v5) node[pos=0.5, right] {$e_5$};
          \draw (v6) to [bend right=25] node[above left] {$e_7$} (v7);
          \draw (v6) to [bend left=25] node[below right] {$e_8$} (v7);
        \end{tikzpicture}
      \end{center}
%    \end{minipage}\hspace*{\stretch{1}}%
%    \begin{minipage}{0.475\linewidth}
%      %Space for matrix
%    \end{minipage}
  \end{ex*}
  \vspace*{\stretch{1}}
  \begin{thmBox*}[Theorem 10.2.1]
    Let $G$ be a graph with connected components $G_1, G_2,\dots, G_k$. If there are $n_i$ vertices in each connected component $G_i$ and these vertices are numbered consecutively, then the adjacency matrix of $G$ has the form
      \[\begin{bmatrix}
        A_1& O& O& \dots& O\\
        O& A_2& O& \dots& O\\
        O& O& A_3& \dots& O\\
        \vdots& \vdots& \vdots& \ddots& \vdots\\
        O& O& O& \dots& A_k
      \end{bmatrix}\]
    where each $A_i$ is the $n_i\times n_i$ adjacency matrix of $G_i$, $\forall i=1,2,\dots,k$, and the $O$'s represent matrices whose entries are all $0$.
  \end{thmBox*}
  \pagebreak

  \begin{defn*}
    Suppose that all elements in matrices $A$ and $B$ are real numbers. If the number of elements $n$ in the $i$th row of $A$ equals the number of elements in the $j$th column of $B$, then the \textbf{scalar product} or \textbf{dot product} of the $i$th row of $A$ and the $j$th column of $B$ is the real number obtained as follows:
      \[
        \begin{bmatrix}
          a_{i1}& a_{i2}& \dots& a_{in}
        \end{bmatrix}
        \begin{bmatrix}
          b_{1j}\\ b_{2j}\\ \vdots\\ b_{nj}
        \end{bmatrix}
        = a_{i1} b_{1j}+ a_{i2}b_{2j} + \dots+ a_{in}b_{nj}.
      \]
  \end{defn*}
  \begin{ex*}
    Multiple the following matrices:
      \[
      \begin{bmatrix}
        -2& -3& 4& 5
      \end{bmatrix}
      \begin{bmatrix}
        -1\\ 0\\ 1\\ -3
      \end{bmatrix}
      \]
  \end{ex*}
  \pagebreak

  \begin{defn*}
    Let $A=(a_{ij})$ be an $m\times k$ matrix, and $B=(b_{ij})$ a $k\times n$ matrix with real entries. The matrix product of $A$ times $B$, denoted $AB$ is the matrix $(c_{ij})$ defined as:
      \[c_{ij}=a_{i1}b_{1j}+a_{i2}b_{2j}+\dots+a_{ik}b_{kj}=\sum_{r=1}^k a_{ir}b_{rj}\]
    for each $i=1,2,\dots,m$ and $j=1,2,\dots,n$.
  \end{defn*}
  \begin{ex*}
    Multiply the following matrices:
    \[
      \begin{bmatrix}
        2&0&-1\\
        -2& 0& 2
      \end{bmatrix}
      \begin{bmatrix}
        5& 5& \\
        -2& -3\\
        3& -4
      \end{bmatrix}
    \]
  \end{ex*}

  \pagebreak
\end{document}
